%!TEX root = ../thesis.tex

\chapter{Future Work}\label{chap:future-work}

% TODO: Follow-up directions, each tied to a limitation established in
% \cref{sec:discussion-disadvantages} or \cref{sec:discussion-threats} rather than
% listed generically. Candidates:
%
%   - Factored action heads with sampled search, lifting the enumerable-branching
%     -factor constraint C5 (\cref{sec:mdp-constraints}) and letting the policy
%     generalize across action factors instead of treating joint actions as unrelated
%     categories (\cref{sec:corpus-action-space}).
%   - Variable-length inputs: a fixed observation window over a variable-length seed
%     decouples "the network needs fixed dimensions" from "seeds must be fixed
%     length", and re-enables insertion and deletion operators.
%   - Coverage representation that scales: edge rather than block coverage, and a
%     collision-aware or learned compression instead of a fixed hashed bucket array
%     (\cref{sec:platform-coverage}).
%   - Reward without shaping: the results depend on hand-designed rewards on the
%     diagnostic targets; intrinsic or novelty-based objectives that need no
%     per-target engineering are the natural next question.
%   - Hybrid deployment: run the agent alongside a conventional fuzzer rather than
%     instead of one, which is how most prior RL fuzzers deal with the throughput gap.
%   - Reliability: collapse is the dominant failure (\cref{chap:failure-modes}), so
%     making runs reliably converge is worth more than making converged runs better.
%   - Broader targets, and a crash-finding rather than coverage-maximizing objective.
